
\chapter{MATLAB Script Analysis: Discrete Math Anonymous Functions}
\author{}
\date{}

\maketitle

\section{Initial Overview}

\textbf{What is the main purpose of this file?}
The main purpose of this file is to define a library or toolkit of compact, one-line anonymous functions in MATLAB. These functions are designed to perform various mathematical and combinatorial operations quickly without the need for standalone function files.

\textbf{What type of analysis or calculation does it perform?}
It performs array size querying, integer divisibility calculations, logical matrix validations for binary relation properties, discrete bijections between sets, random integer partitioning, and conversions between different representations of graphs/permutations (such as signatures, connected components, and cycles).

\textbf{What is the mathematical/theoretical context?}
The theoretical context is heavily rooted in Discrete Mathematics and Combinatorics. Specifically, it deals with:
\begin{itemize}
    \item \textbf{Binary Relations \& Posets:} Checking for reflexivity, transitivity, symmetry, and antisymmetry using adjacency matrices.
    \item \textbf{Number Theory:} Finding divisors and building Hasse diagrams for the divisibility poset.
    \item \textbf{Set Theory \& Morphisms:} Calculating bijections between functional sets ($B^X$) and composing endomaps.
    \item \textbf{Graph Theory \& Permutation Groups:} Mapping integer partitions (signatures) to connected components and disjoint cycles.
\end{itemize}

\section{Step-by-Step Code Analysis}

\subsection{1. Section/Function Title: Matrix and Vector Sizing Utilities}
\textbf{2. Explanation:} This block contains basic helper functions to query the dimensions of arrays. \texttt{nlin} returns the number of rows, and \texttt{ncol} returns the number of columns. The functions \texttt{idn} and \texttt{id} generate vertical column vectors representing a sequence of indices, up to a specified integer or the total number of elements in an array, respectively.

\textbf{3. Code Snippet:}
\begin{lstlisting}
nlin=@(x) size(x,1);
ncol=@(x) size(x,2);
idn=@(x) (1:x)';
id=@(x) (1:numel(x))';
\end{lstlisting}

\subsection{1. Section/Function Title: Divisors and Hasse Diagram Relation}
\textbf{2. Explanation:} The \texttt{div} function calculates all divisors of a given integer $n$ by finding where the modulo operation yields zero. Following this, the code explicitly calculates the divisors for $n=1000$ and builds a boolean adjacency matrix \texttt{R}. This matrix represents the divisibility relation $a|b$ among the divisors, acting as the foundation for a Hasse diagram. 
\textbf{3. Code Snippet:}
\begin{lstlisting}
% collection of divisors of a given number
div=@(n) find(mod(n,1:n)==0);
% and respective Hasse diagram
u=div(1000); [x,y]=ndgrid(id(u));
R=(mod(u(y),u(x))==0); 
\end{lstlisting}

\subsection{1. Section/Function Title: Properties of Binary Relations}
\textbf{2. Explanation:} These anonymous functions evaluate a square adjacency matrix \texttt{R} representing a binary relation to determine its mathematical properties. 
\begin{itemize}
    \item \texttt{isreflexive} checks if the main diagonal is all ones.
    \item \texttt{istransitive} checks if the matrix squared (representing two-step paths) is a subset of the original matrix.
    \item \texttt{issymmetric} checks if the matrix is equal to its transpose.
    \item \texttt{isantisymmetric} ensures that if $aRb$ and $bRa$, then $a=b$ (i.e., no bidirectional links outside the main diagonal).
\end{itemize}

\textbf{3. Code Snippet:}
\begin{lstlisting}
% properties of relations
isreflexive=@(R) all(diag(R));
istransitive=@(R) all(all(logical(R*R)<=R));
issymmetric=@(R) isequal(R,R');
isantisymmetric=@(R) all(all((R&R')<=eye(size(R))));
\end{lstlisting}

\subsection{1. Section/Function Title: Endomaps and Functional Bijections}
\textbf{2. Explanation:} This section deals with maps between finite sets. \texttt{composerows} executes a fast matrix-based composition of functions. \texttt{s2i} (sequence to index) and \texttt{i2s} (index to sequence) establish a bijection between sequences (morphisms) and integer indices, effectively encoding and decoding elements of the set $B^X$. Finally, \texttt{circ} composes two morphisms represented purely by their integer indices.

\textbf{3. Code Snippet:}
\begin{lstlisting}
% the composition of endomaps as if they were row vectors
composerows=@(A,J) A(sub2ind(size(J),repmat((1:size(J,1))',1,size(J,2)),J));
% the bijection between A and B^X
s2i=@(nB,s) (s-1)*(nB.^((1:size(s,2))'-1))+1; 
i2s=@(i,nB,nX) mod(floor((i(:)-1)./nB.^((1:nX)-1)),nB)+1; 
circ=@(i,j,nB,nX,nZ) s2i(composerows(i2s(i,nB,nX),i2s(j,nX,nZ)),nB); 
\end{lstlisting}

\subsection{1. Section/Function Title: Integer Partitioning}
\textbf{2. Explanation:} The \texttt{partition} function randomly divides an integer $n$ into $k$ distinct, strictly positive integer parts. It achieves this by generating $k-1$ random cut points in the interval $[1, n-1]$ and taking the differences between these sorted cut points, mimicking the "stars and bars" combinatorial method.

\textbf{3. Code Snippet:}
\begin{lstlisting}
% random partition 
partition=@(n,k) diff([0 sort(randperm(n-1,k-1)) n]); 
\end{lstlisting}

\subsection{1. Section/Function Title: Signatures, Connected Components, and Cycles}
\textbf{2. Explanation:} These functions transform graph/permutation data representations. A "signature" here acts as an array of counts or sizes. 
\begin{itemize}
    \item \texttt{siginterp} rescales a signature to sum to a new total $n$.
    \item \texttt{sig2conn} and \texttt{conn2sig} map between a signature (e.g., sizes of connected components) and a flat array assigning each element to a component ID.     \item \texttt{sig2cycles} constructs a canonical permutation mapping from cycle sizes, and \texttt{cycles2sig} attempts to invert this process (relying on an undefined external \texttt{orbits} function).
\end{itemize}

\textbf{3. Code Snippet:}
\begin{lstlisting}
% signature interpolation
siginterp=@(sig,n) diff([0, round(cumsum(sig)/sum(sig)*n)]);

% signature to connected components and back
sig2conn=@(sig) repelem(1:numel(sig),sig)
conn2sig=@(conn) full(sparse(conn,1,1))' 

% signature to cycles and back
sig2cycles=@(sig) 1+idn(sum(sig))+sparse(cumsum(sig),1,-sig);
cycles2sig=@(phi) conn2sig(orbits(phi))
\end{lstlisting}

\section{Expected Outputs and Visuals}

When this script is executed, the following outputs and data structures are generated:

\begin{itemize}
    \item \textbf{Console Output:} The script immediately outputs the string literal: \texttt{List of useful anonymous functions} to the command window.
    \item \textbf{The Relation Matrix ($R$):} The script quietly calculates a matrix $R$ for the divisors of 1000. The divisors of 1000 are: $1, 2, 4, 5, 8, 10, 20, 25, 40, 50, 100, 125, 200, 250, 500, 1000$ (16 divisors in total). 
    \begin{itemize}
        \item $R$ will be a $16 \times 16$ logical matrix.
        \item The $x$-axis (columns) and $y$-axis (rows) represent the indices of these sorted divisors.
        \item A value of \texttt{1} (True) at position $(y, x)$ indicates that the divisor at index $x$ perfectly divides the divisor at index $y$ (i.e., $u(x) \mid u(y)$).
        \item \textbf{Visual Translation:} If visualized using MATLAB's \texttt{spy(R)} or plotted as a directed graph using \texttt{digraph(R)}, it visually manifests as a partially ordered set (poset). The main diagonal would be fully populated (reflexive), and the lower triangle would contain scattered points mapping the branching hierarchy of multiples, essentially representing a flattened Hasse diagram of the prime factorization $2^3 \times 5^3$.
    \end{itemize}
\end{itemize}
